\section{操作手册与工具参考}

\subsection{Reward Learning Action Timeline}

\begin{table}[H]
\centering
\begin{tabular}{p{0.25\textwidth} p{0.45\textwidth} p{0.2\textwidth}}
\toprule
阶段 & 主要活动 & 负责人 \\
\midrule
数据准备 & 采集 positive/negative 示例、偏好比较 & Data engineer + labelers \\
Reward training & 训练 classifier/RLM、监控 loss 曲线 & ML engineer \\
Evaluation & 自动指标、human eval、safety probe & Ops + safety reviewer \\
Deployment & KL penalty tuning、logging + rollback & DevOps \\
\bottomrule
\end{tabular}
\caption{Reward learning 的操作时间线}
\end{table}

\subsection{工具与命令}

\begin{itemize}
    \item 使用 `tensorboard` 监控 reward model 的 logits + loss
    \item 将 `reward_score` 写入 experiment tracker（Weights & Biases / Neptune）
    \item RLHF pipeline 推荐用 `trl` + `DPO` + `PPO` mix
\end{itemize}

\begin{knowledgebox}{资源定位}
\textbf{Reward learning 资料}集中在：OpenAI blog、Anthropic interpretability notes、Stanford RL lab releases。保持订阅这些更新可以及时调整 payoff curve。
\end{knowledgebox}

\begin{figure}[H]
\centering
\includegraphics[width=0.92\textwidth,page=7]{08_cs224r_reward_learning_2025.pdf}
\caption{Reward trace dashboard 例示（lecture slide, page 7）}
\end{figure}

\subsection{本章小结}

操作手册强调每一步需要负责人并行推进，尤其是 reward training 与 evaluation 不能脱钩。

